% =====================================================================
%  CHAPTER 7: 儿童少年健康状况
% =====================================================================
\chapter{儿童少年健康状况}
\setcounter{chapter}{7}
\chapterintro{
掌握健康的三维度（生物-心理-社会医学模式）；掌握儿童少年健康评价指标体系；掌握儿童少年常见健康问题和疾病谱变化；熟悉儿童少年健康问题的流行病学现状（超重肥胖、营养不良、心理健康等）；理解健康促进在儿童少年群体中的应用。本章"儿童期和青春期的健康特点"为教学难点（*）。
}

% ----- 7.1 健康概念与评价 -----
\section{健康的概念与评价指标体系}

\subsection{健康的概念}

\begin{defbox}
\textbf{健康（Health）}：WHO定义——健康不仅是躯体没有疾病，还须具备心理健康、社会适应良好和道德健康。是身体、心理、社会适应和道德上的完美状态。

\textbf{亚健康（Sub-Health）}：处于健康与疾病之间的一种状态，个体在身体、心理和社会适应上出现不适感但尚未达到疾病诊断标准。反映某些症状但传统临床检验可无明显异常。
\end{defbox}

\subsection{基于生物-心理-社会医学模式的健康维度}

\begin{keybox}
\textbf{健康的四个维度及评价指标：}

\begin{enumerate}[leftmargin=*]
    \item \imp{身体健康指标}：生长发育水平、生理功能、运动能力、日常生活活动能力等
    \item \imp{心理健康指标}：情绪/情感症状、行为表现、认知功能、性格和个性、主观幸福感等
    \item \imp{社会适应指标}：人际关系、学习适应能力、日常生活适应能力、社会参与等
    \item \imp{道德健康指标}：道德认知、道德情感、道德意志、道德行为习惯等
\end{enumerate}
\end{keybox}

\subsection{疾病发展过程的健康维度}

\begin{defbox}
\textbf{基于疾病发展过程的健康评价指标：}

\begin{enumerate}[leftmargin=*]
    \item \imp{健康状态}：无疾病、身心社会适应完好
    \item \imp{亚健康状态}：介于健康与疾病之间，有不适感但无客观检查异常
    \item \imp{发病率、患病率和死亡率}：
    \begin{itemize}
        \item 发病率——反映健康风险的动态变化及防控措施效果
        \item 患病率——反映疾病在人群中的普遍性和慢性病的管理负担
        \item 死亡率——最直接揭示健康问题的严重程度及对儿童生命的威胁
    \end{itemize}
    \item \imp{医疗卫生服务利用}：预防接种、健康检查、医疗服务利用、医疗费用等
    \item \imp{疾病负担指标}：DALY（伤残调整寿命年）、因病缺课天数等
\end{enumerate}
\end{defbox}

% ----- 7.2 儿童少年健康现况 -----
\section{儿童少年健康问题的流行病学现况}

\subsection{全球及中国儿童健康挑战}

\begin{keybox}
\textbf{（一）营养与生长发育挑战依然严峻：}

\begin{enumerate}[leftmargin=*]
    \item \imp{超重肥胖发生率持续上升}
    \begin{itemize}
        \item 全球：2022年约3700万5岁以下儿童超重，增长主要在发展中国家
        \item 中国：2020年6岁以下儿童超重率达10.4\%，6--17岁达19\%，预计2027年农村儿童超重肥胖患病率将接近城市水平
    \end{itemize}

    \item \imp{营养不良问题仍然存在}
    \begin{itemize}
        \item 全球：2020年22\%的5岁以下儿童生长迟缓，6.7\%消瘦
        \item 中国：6岁以下生长迟缓率4.8\%（农村高于城市），6--17岁消瘦率2.2\%，贫血率6.1\%
    \end{itemize}

    \item \imp{心血管代谢危险因素需重点关注}
    \begin{itemize}
        \item 心血管和代谢疾病风险因素普遍化
        \item 膳食、睡眠和体力活动达标率低
    \end{itemize}
\end{enumerate}

\textbf{（二）心理健康问题日益突出（难点*）：}

随着社会经济发展、城市化加速、生活方式改变、学业压力增大，儿童少年群体的心理健康问题越来越突出。

\begin{itemize}[leftmargin=*]
    \item 全球：约20\%的10--19岁青少年存在心理健康问题（UNICEF/WHO报告）
    \item 中国：约14.8\%的青少年存在不同程度的心理健康风险，其中4.0\%为重度风险群体，10.8\%为轻度风险群体
\end{itemize}

\textbf{（三）疾病谱的变化：}

儿童少年疾病谱发生显著转变：
\begin{itemize}[leftmargin=*]
    \item 传染病和感染性疾病发病率和死亡率显著下降
    \item 慢性非传染性疾病（肥胖、糖尿病、高血压、哮喘等）成为日益重要的健康问题
    \item 心理健康问题、意外伤害、物质滥用（吸烟、饮酒、网瘾等）构成新的健康威胁
\end{itemize}
\end{keybox}

% ----- 7.3 儿童期和青春期的健康特点 -----
\section{儿童期和青春期的健康特点（教学难点*）}

\subsection{儿童期健康特点}

\begin{defbox}
\textbf{儿童期（学龄期）健康特点：}

\begin{enumerate}[leftmargin=*]
    \item \imp{各系统功能逐渐完善但尚不成熟}——免疫功能逐渐增强但仍有感染风险；视觉系统处于敏感发育期（近视防控关键期）
    \item \imp{生长发育稳定但个体差异明显}
    \item \imp{常见健康问题}：近视、龋齿、营养不良/超重肥胖、脊柱弯曲异常、感染性疾病
    \item \imp{行为习惯形成期}——饮食、运动、卫生、学习习惯在此时期形成
    \item \imp{意外伤害高发}——交通事故、溺水、跌落等
\end{enumerate}
\end{defbox}

\subsection{青春期健康特点}

\begin{defbox}
\textbf{青春期健康特点（难点*）：}

\begin{enumerate}[leftmargin=*]
    \item \imp{身心发育的巨大变化}——第二生长突增、性发育启动并完成，带来新的健康需求和问题
    \item \imp{心理行为发育的关键窗口}——情绪波动、冒险行为增多、自我认同探索
    \item \imp{生殖健康问题}——青春期妊娠、性传播疾病、月经异常
    \item \imp{慢性病危险因素累积}——不良饮食、缺乏运动、吸烟饮酒等危险行为多在青春期启动
    \item \imp{心理健康高风险}——抑郁、焦虑、自伤自残行为发生率在青春期显著升高
    \item \imp{体象关注与饮食失调}——对身体外貌过度关注，可能导致节食、暴食等饮食行为问题
\end{enumerate}
\end{defbox}

% ----- 7.4 健康促进 -----
\section{健康促进在儿童少年群体中的应用}

\begin{defbox}
\textbf{健康促进（Health Promotion）}：运用行政或组织手段，广泛动员和协调社会各相关部门以及社区、家庭和个人，使其履行各自对健康的责任，共同维护和促进健康的一种社会行为和社会战略。

\textbf{儿童少年健康促进的基本原则：}
\begin{enumerate}[leftmargin=*]
    \item \imp{全人群策略与高危策略结合}——面向全体儿童同时关注高危人群
    \item \imp{多部门协作}——卫生、教育、体育、食品等多部门共同参与
    \item \imp{学校为基础}——学校是健康促进最重要的场所
    \item \imp{家庭参与}——家庭是儿童健康的第一道防线
    \item \imp{全生命周期视角}——关注儿童健康对终身健康的影响
    \item \imp{健康融入所有政策（HiAP）}——所有公共政策都应考虑对儿童健康的影响
\end{enumerate}
\end{defbox}

% ----- 复习题 -----
\section{复习题}

\begin{enumerate}[leftmargin=*, itemsep=8pt]
    \item \textbf{【名词解释】}健康（WHO定义）；亚健康；健康促进
    \item \textbf{【名词解释】}DALY；疾病谱
    \item \textbf{【简答题】}简述基于生物-心理-社会医学模式的健康四维度评价指标。
    \item \textbf{【简答题】}试述当前中国儿童少年面临的主要健康问题和挑战。
    \item \textbf{【简答题】}简述儿童期与青春期的健康特点。
    \item \textbf{【简答题】}简述儿童少年健康促进的基本原则。
    \item \textbf{【简答题】}试分析儿童少年疾病谱变化的主要趋势及原因。
\end{enumerate}

\subsection*{参考答案要点}

\begin{enumerate}[leftmargin=*, itemsep=5pt]
    \item \textbf{健康}：身体、心理、社会适应和道德上的完美状态，不仅是没有疾病。\textbf{亚健康}：健康与疾病之间的一种状态，有不适感但无客观检查异常。\textbf{健康促进}：运用行政组织手段动员社会力量共同维护和促进健康的社会战略。

    \item \textbf{DALY}：伤残调整寿命年，量化疾病负担的综合指标（早死所致生命损失年+伤残所致健康生命损失年）。\textbf{疾病谱}：人群疾病构成的变化格局。

    \item 身体健康（生长发育、生理功能、运动能力）、心理健康（情绪、行为、认知、幸福感）、社会适应（人际关系、学习适应、社会参与）、道德健康（道德认知、情感、意志、行为习惯）。

    \item (1)营养不良与超重肥胖双重负担；(2)心理健康问题日益突出；(3)慢性非传染性疾病危险因素普遍化；(4)意外伤害是主要死因之一；(5)疾病谱从传染病为主转向慢性病和心理健康问题为主。

    \item 儿童期：系统功能完善中但未成熟、视觉敏感期、常见病（近视、龋齿）、行为习惯形成期、意外伤害高发。青春期：身心剧变、心理行为关键窗口、生殖健康问题、慢性病危险因素积累、心理健康高风险、体象关注。

    \item (1)全人群与高危策略结合；(2)多部门协作（卫生、教育、体育等）；(3)学校为基础；(4)家庭参与；(5)全生命周期视角；(6)健康融入所有政策。

    \item 趋势：传染病和感染病降低→慢性非传染性疾病上升→心理健康问题和新发问题（意外伤害、物质滥用、网瘾等）增多。原因：经济社会发展、城市化、生活方式改变、学业压力增大、预防接种推广等综合作用。
\end{enumerate}

\newpage
